% ============================================================
\chapter{Complexes de Vietoris-Rips, \v{C}ech et Alpha}
\label{ch:complexes}
% ============================================================

L'homologie persistante a besoin d'une \emph{filtration} de complexes simpliciaux. Mais comment construire un complexe \`a partir d'un nuage de points~? Trois constructions classiques dominent. Le complexe de \emph{\v{C}ech} ajoute un simplexe d\`es que les boules de rayon $r$ centr\'ees sur ses sommets ont une intersection commune --- il capture exactement la topologie de la r\'eunion des boules (th\'eor\`eme du nerf). Le complexe de \emph{Vietoris-Rips} est plus simple~: un simplexe est pr\'esent d\`es que toutes les paires de sommets sont \`a distance $\leq 2r$ --- il est bien plus facile \`a calculer, mais approxime seulement le \v{C}ech. Le complexe \emph{Alpha}, d\^u \`a Edelsbrunner, restreint le \v{C}ech aux cellules de Vorono\"i, produisant un complexe de taille lin\'eaire en le nombre de points.

\begin{intuition}
La construction d'un complexe simplicial \`a partir d'un nuage de points
est l'\'etape fondamentale de la TDA. Trois constructions dominent~:
Vietoris--Rips (simple mais gros), \v{C}ech (exact mais co\^uteux),
et Alpha (petit et exact en basse dimension).
\end{intuition}

% ============================================================
\section{Complexe de Vietoris--Rips}
% ============================================================

\begin{definition}[Complexe de Vietoris--Rips]
Soit $X = \{x_1, \ldots, x_n\}$ un ensemble fini dans un espace
m\'etrique $(M, d)$. Le \emph{complexe de Vietoris--Rips} au
param\`etre $r \geq 0$ est~:
\[
  \mathrm{VR}(X, r) = \big\{
    \sigma \subseteq X : d(x_i, x_j) \leq r
    \;\text{pour tous}\; x_i, x_j \in \sigma
  \big\}.
\]
\end{definition}

\begin{remarque}
Le complexe de Vietoris--Rips est le \emph{complexe de cliques}
(ou \emph{flag complex}) du graphe de proximit\'e $G_r = (X, E_r)$
o\`u $\{x_i, x_j\} \in E_r$ si et seulement si $d(x_i, x_j) \leq r$.
Autrement dit, $\sigma \in \mathrm{VR}(X, r)$ si et seulement si toutes
les paires de sommets de $\sigma$ sont reli\'ees dans $G_r$.
\end{remarque}

\begin{proposition}[Propri\'et\'es]
\begin{enumerate}
  \item $\mathrm{VR}(X, r)$ est enti\`erement d\'etermin\'e par son
        $1$-squelette (le graphe $G_r$).
  \item $r \leq s \implies \mathrm{VR}(X, r) \subseteq \mathrm{VR}(X, s)$.
  \item La dimension maximale de $\mathrm{VR}(X, r)$ peut atteindre
        $n - 1$.
\end{enumerate}
\end{proposition}

\begin{attention}[title=\textbf{Explosion combinatoire}]
Pour $n$ points, le complexe de Vietoris--Rips peut contenir
jusqu'\`a $2^n - 1$ simplexes. En pratique, on tronque la dimension
maximale (typiquement $\leq 3$).
\end{attention}

% ============================================================
\section{Complexe de \v{C}ech}
% ============================================================

\begin{definition}[Complexe de \v{C}ech]
Le \emph{complexe de \v{C}ech} au param\`etre $r$ est~:
\[
  \check{C}(X, r) = \big\{
    \sigma \subseteq X : \bigcap_{x_i \in \sigma} B(x_i, r) \neq \emptyset
  \big\}
\]
o\`u $B(x_i, r) = \{y \in M : d(x_i, y) \leq r\}$ est la boule
ferm\'ee de centre $x_i$ et de rayon $r$.
\end{definition}

\begin{theoreme}[Nerf --- Borsuk, 1948]
Soit $\mathcal{U} = \{U_i\}_{i \in I}$ un recouvrement fini d'un
espace topologique $X$ par des ouverts convexes. Alors le \emph{nerf}
de $\mathcal{U}$~:
\[
  \mathcal{N}(\mathcal{U}) = \big\{
    \sigma \subseteq I : \bigcap_{i \in \sigma} U_i \neq \emptyset
  \big\}
\]
a le m\^eme type d'homotopie que $\bigcup_{i \in I} U_i$.
\end{theoreme}

\begin{corollaire}
Dans $\R^d$, $\check{C}(X, r)$ a le m\^eme type d'homotopie que
$\bigcup_{i=1}^n B(x_i, r)$.
\end{corollaire}

\begin{proposition}[Relation Rips--\v{C}ech]
Pour tout nuage de points $X \subset \R^d$ et tout $r > 0$~:
\[
  \check{C}(X, r) \subseteq \mathrm{VR}(X, 2r) \subseteq
  \check{C}\left(X, \frac{2r}{\sqrt{2}} \cdot \sqrt{\frac{d}{d+1}}\right).
\]
En particulier~: $\check{C}(X, r) \subseteq \mathrm{VR}(X, 2r)$.
\end{proposition}

% ============================================================
\section{Complexe Alpha}
% ============================================================

\begin{definition}[Diagramme de Vorono\"{\i}]
Le \emph{diagramme de Vorono\"{\i}} de $X \subset \R^d$ est la
partition de $\R^d$ en cellules~:
\[
  V_i = \{y \in \R^d : \norm{y - x_i} \leq \norm{y - x_j}
  \;\text{pour tout}\; j\}.
\]
\end{definition}

\begin{definition}[Triangulation de Delaunay]
La \emph{triangulation de Delaunay} $\mathrm{Del}(X)$ est le complexe
simplicial dual du diagramme de Vorono\"{\i}~:
$\sigma = \{x_{i_0}, \ldots, x_{i_p}\} \in \mathrm{Del}(X)$ si et
seulement si $\bigcap_{k=0}^p V_{i_k} \neq \emptyset$.
\end{definition}

\begin{definition}[Complexe Alpha]
Le \emph{complexe Alpha} au param\`etre $r$ est~:
\[
  \mathrm{Alpha}(X, r) = \big\{
    \sigma \in \mathrm{Del}(X) :
    \bigcap_{x_i \in \sigma} (B(x_i, r) \cap V_i) \neq \emptyset
  \big\}.
\]
C'est le sous-complexe de la triangulation de Delaunay form\'e des
simplexes dont les boules restreintes aux cellules de Vorono\"{\i}
s'intersectent.
\end{definition}

\begin{theoreme}[Edelsbrunner, 1995]
Le complexe Alpha v\'erifie~:
\begin{enumerate}
  \item $\mathrm{Alpha}(X, r)$ a le m\^eme type d'homotopie que
        $\bigcup_{i=1}^n B(x_i, r)$ (et donc que $\check{C}(X, r)$).
  \item $\dim(\mathrm{Alpha}(X, r)) \leq d$ (la dimension ambiante).
  \item Le nombre total de simplexes est $O(n^{\lceil d/2 \rceil})$.
\end{enumerate}
\end{theoreme}

\begin{formules}[title=\textbf{Comparaison des tailles}]
Pour $n$ points dans $\R^d$~:
\begin{center}
\begin{tabular}{lcc}
\toprule
\textbf{Complexe} & \textbf{Dimension max} & \textbf{Nb simplexes} \\
\midrule
Vietoris--Rips & $n - 1$ & $O(2^n)$ \\
\v{C}ech & $n - 1$ & $O(2^n)$ \\
Alpha & $d$ & $O(n^{\lceil d/2 \rceil})$ \\
\bottomrule
\end{tabular}
\end{center}
\end{formules}

% ============================================================
\section{Construction en pratique}
% ============================================================

\begin{codeblock}[title=\textbf{Complexe de Rips avec GUDHI}]
\begin{minted}{python}
import gudhi
import numpy as np

# Nuage de points
np.random.seed(42)
X = np.random.randn(50, 3)

# Complexe de Rips
rips = gudhi.RipsComplex(points=X, max_edge_length=2.0)
st = rips.create_simplex_tree(max_dimension=3)

print(f"Nombre de simplexes: {st.num_simplices()}")
print(f"Nombre de sommets: {st.num_vertices()}")
print(f"Dimension: {st.dimension()}")

# Persistance
pers = st.persistence()
gudhi.plot_persistence_diagram(pers)
\end{minted}
\end{codeblock}

\begin{codeblock}[title=\textbf{Complexe Alpha avec GUDHI}]
\begin{minted}{python}
import gudhi
import numpy as np

# Nuage de points en 2D
np.random.seed(42)
theta = np.linspace(0, 2*np.pi, 80, endpoint=False)
X = np.column_stack([np.cos(theta), np.sin(theta)])
X += 0.1 * np.random.randn(*X.shape)

# Complexe Alpha
alpha = gudhi.AlphaComplex(points=X)
st = alpha.create_simplex_tree()

print(f"Nombre de simplexes: {st.num_simplices()}")
print(f"Dimension: {st.dimension()}")

# Persistance
pers = st.persistence()
gudhi.plot_persistence_diagram(pers)
\end{minted}
\end{codeblock}

\begin{codeblock}[title=\textbf{Complexe de \v{C}ech avec GUDHI}]
\begin{minted}{python}
import gudhi
import numpy as np

# Nuage de points en 2D
X = np.random.randn(30, 2)

# Le complexe de Čech n'a pas de classe dédiée dans GUDHI
# mais on peut utiliser le complexe de Rips comme approximation
# ou utiliser miniball pour vérifier le critère de Čech

# Alternative: utiliser la filtration Alpha qui est
# homotopiquement équivalente au Čech
alpha = gudhi.AlphaComplex(points=X)
st = alpha.create_simplex_tree()

# Les valeurs de filtration Alpha sont les carrés des rayons
# de la plus petite boule englobante de chaque simplexe
for simplex, filt in st.get_filtration():
    if len(simplex) <= 3:
        print(f"  {simplex}: rayon^2 = {filt:.4f}")
    if filt > 0.5:
        break
\end{minted}
\end{codeblock}

% ============================================================
\section{Complexe de Rips pondéré et Rips spars}
% ============================================================

\begin{definition}[Sparse Rips --- Sheehy, 2013]
Le \emph{Sparse Rips complex} est une approximation du complexe de
Vietoris--Rips qui contient $O(n)$ simplexes dans chaque dimension, au
lieu de $O(2^n)$. Pour un param\`etre $\varepsilon > 0$, il produit
une $(1+\varepsilon)$-approximation des diagrammes de persistance.
\end{definition}

\begin{codeblock}[title=\textbf{Sparse Rips avec GUDHI}]
\begin{minted}{python}
import gudhi
import numpy as np

X = np.random.randn(500, 3)

# Sparse Rips (epsilon = 0.3)
sparse_rips = gudhi.RipsComplex(
    points=X, max_edge_length=2.0, sparse=0.3)
st = sparse_rips.create_simplex_tree(max_dimension=2)

print(f"Sparse Rips: {st.num_simplices()} simplexes")

# Comparaison avec Rips standard
rips = gudhi.RipsComplex(points=X[:100], max_edge_length=2.0)
st_full = rips.create_simplex_tree(max_dimension=2)
print(f"Full Rips (100 pts): {st_full.num_simplices()} simplexes")
\end{minted}
\end{codeblock}

% ============================================================
\section{Complexes cubiques}
% ============================================================

\begin{definition}[Complexe cubique]
Un \emph{complexe cubique} est l'analogue d'un complexe simplicial
o\`u les briques \'el\'ementaires sont des cubes (produits d'intervalles)
au lieu de simplexes. Il est particuli\`erement adapt\'e aux donn\'ees
sur grille (images, volumes 3D).
\end{definition}

\begin{codeblock}[title=\textbf{Homologie persistante d'une image}]
\begin{minted}{python}
import gudhi
import numpy as np

# Créer une image simple (anneau)
x = np.linspace(-2, 2, 50)
y = np.linspace(-2, 2, 50)
XX, YY = np.meshgrid(x, y)
image = np.sqrt(XX**2 + YY**2)  # distance à l'origine

# Complexe cubique
cc = gudhi.CubicalComplex(
    top_dimensional_cells=image.flatten(),
    dimensions=image.shape
)

# Persistance
pers = cc.persistence()
print("Persistance de l'image:")
for dim, (b, d) in pers:
    if d - b > 0.5:  # filtrer le bruit
        print(f"  H_{dim}: ({b:.2f}, {d:.2f}), "
              f"pers={d-b:.2f}")
\end{minted}
\end{codeblock}

% ============================================================
\section{Choix du complexe en pratique}
% ============================================================

\begin{algorithme}[title=\textbf{Guide de choix du complexe}]
\begin{itemize}
  \item \textbf{Donn\'ees en basse dimension} ($d \leq 3$)~:
        utiliser le complexe \textbf{Alpha}. Il est petit, exact, et
        rapide \`a construire.
  \item \textbf{Donn\'ees en haute dimension ou m\'etriques}~:
        utiliser \textbf{Vietoris--Rips} (avec \texttt{ripser} pour la
        vitesse). Tronquer la dimension \`a $2$ ou $3$.
  \item \textbf{Donn\'ees sur grille} (images, volumes)~: utiliser un
        complexe \textbf{cubique}.
  \item \textbf{Grand nombre de points}~: utiliser \textbf{Sparse Rips}
        ou sous-\'echantillonner.
\end{itemize}
\end{algorithme}

% ============================================================
\section{Exercices}
% ============================================================

\begin{exercice}
Pour $4$ points aux sommets d'un carr\'e unit\'e, construire \`a la
main les complexes $\mathrm{VR}(X, r)$ et $\check{C}(X, r)$ pour
$r = 0.5, 1.0, 1.5$. V\'erifier l'inclusion
$\check{C}(X, r) \subseteq \mathrm{VR}(X, 2r)$.
\end{exercice}

\begin{exercice}
Construire le diagramme de Vorono\"{\i} et la triangulation de Delaunay
pour $5$ points al\'eatoires dans $\R^2$ (\'eventuellement avec
\texttt{scipy.spatial.Delaunay}).
\end{exercice}

\begin{exercice}
Comparer les temps de calcul des complexes Alpha et Rips pour des
nuages de $100$ \`a $1000$ points en dimensions $2$, $3$ et $10$.
\end{exercice}

\begin{exercice}
Montrer que le complexe de Vietoris--Rips d'un ensemble de $n+1$
points en position g\'en\'erale dans $\R^n$, pour $r$ suffisamment
grand, est le simplexe plein $\Delta^n$.
\end{exercice}

\begin{exercice}
Calculer l'homologie persistante d'une image binaire de votre choix
en utilisant un complexe cubique avec GUDHI.
\end{exercice}
